\chapter{Introdução}

\section{Contextualização}

Nas últimas décadas, a relação entre Estado e cidadão passou por uma mudança
estrutural: atividades que antes dependiam de balcões presenciais, deslocamentos
e atendimento físico foram progressivamente incorporadas a portais, aplicações e
fluxos digitais. A plataforma GOV.BR simboliza esse movimento ao concentrar a
oferta digital de serviços federais e ao reforçar a ideia de que governo digital
não é apenas troca de canal, mas reorganização da forma como o cidadão encontra,
solicita, acompanha e avalia serviços públicos.

Os serviços públicos compõem uma das formas mais diretas dessa relação. Solicitar
um documento, consultar uma situação cadastral, acompanhar um benefício,
registrar uma demanda ou acessar uma política pública são situações em que a
experiência do usuário influencia a qualidade percebida da administração
pública. A digitalização ampliou o alcance e a velocidade desses atendimentos,
mas também deslocou parte das dificuldades para a própria navegação: localizar o
serviço correto, entender a linguagem institucional, preencher etapas
intermediárias e perceber se a solicitação realmente avançou.

A escala dessa transição aparece de forma clara em experiências recentes de
participação digital. O Plano Plurianual Participativo 2024-2027, por exemplo,
mobilizou mais de 1,5 milhão de participantes e superou 4 milhões de acessos na
plataforma Brasil Participativo, demonstrando a capacidade de mobilização do
ambiente digital. Ao mesmo tempo, a experiência registrou barreiras como
problemas de conectividade, dificuldade de uso, fluxos de cadastro complexos e
necessidade de melhor retorno ao cidadão sobre sua contribuição
\cite{brasilMgi2024ppaParticipativo}. Esse contraste reforça a importância de
avaliar não apenas se o serviço existe em formato digital, mas como ele é
percebido e utilizado durante a jornada.

No Brasil, a avaliação dos serviços públicos possui base legal e gerencial. A
Lei n\textsuperscript{o} 13.460/2017 trata da participação, proteção e defesa dos
usuários de serviços públicos e prevê que os resultados de avaliação subsidiem a
reorientação e o ajuste da prestação estatal \cite{brasil2017lei13460}. O
Decreto n\textsuperscript{o} 9.094/2017 institui a Carta de Serviços ao Usuário
e reforça a simplificação do atendimento \cite{brasil2017decreto9094}. A Lei
n\textsuperscript{o} 14.129/2021 associa governo digital, eficiência pública,
desburocratização e participação do cidadão \cite{brasil2021lei14129}. O Decreto
n\textsuperscript{o} 11.260/2022, ao tratar da Estratégia Nacional de Governo
Digital, reforça esse movimento de organização e monitoramento da transformação
digital na administração pública federal \cite{brasil2022decreto11260}.

A avaliação da experiência digital, porém, não depende apenas de perguntar ao
usuário se ele ficou satisfeito ao final do atendimento. Em serviços digitais, a
dificuldade pode aparecer durante a navegação: excesso de cliques, permanência
por tempo acima do esperado, retorno repetido a páginas anteriores, busca por
serviços prioritários e abandono antes da conclusão. Métricas de uso, quando
coletadas com finalidade clara e com cuidado em relação à privacidade, ajudam a
identificar gargalos que um formulário isolado dificilmente revela
\cite{deloneMclean2003,govuk2021completionRate,brasil2018lgpd}.

A proposta deste TCC é criar um avaliador de serviços públicos para ambientes
digitais que una dois tipos de evidência. O primeiro tipo é comportamental: a API
do protótipo registra cliques, página acessada, serviço associado, origem e tempo
de navegação, calculando sinais de dificuldade a partir de faixas esperadas. O
segundo tipo é declarativo: quando o usuário apresenta sinais de dificuldade ou
permanece mais tempo que a média, o sistema pode abrir uma página simples de
feedback com avaliação por estrelas, características do problema e uma pergunta
de NPS. Para serviços marcados como prioritários ou de alto tráfego, como Imposto
de Renda, o protótipo pode sugerir um redirecionamento ao caminho correto antes
de solicitar o feedback.

A ideia se inspira na simplicidade de avaliações presentes em plataformas
digitais conhecidas pelo público, mas não pretende copiar sua lógica comercial.
Em serviços públicos, a avaliação precisa apoiar participação social,
transparência, melhoria contínua e priorização de correções. Por isso, o
avaliador proposto combina interface simples para o cidadão, registro silencioso
de sinais de dificuldade e visualização administrativa de alertas, sessões
críticas, notas, NPS e pontos de melhoria. A versão desenvolvida neste TCC foi
definida por meio de Lean Inception, usada para delimitar o MVP, suas personas,
jornadas, funcionalidades e critérios de validação
\cite{caroli2018leanInception,fowler2022leanInception}.

\section{Problema de Pesquisa}

Apesar do avanço quantitativo na oferta digital de serviços públicos, a avaliação
da experiência ainda tende a ocorrer tarde, fora do contexto em que a dificuldade
aparece. Formulários exibidos apenas ao final da jornada, pesquisas esporádicas e
reclamações registradas depois do problema ajudam a compor diagnósticos, mas nem
sempre capturam o momento em que o usuário se perde, repete cliques, permanece
tempo excessivo em uma etapa ou abandona o fluxo. Essa distância entre a
dificuldade vivida e a leitura institucional do problema reduz a capacidade de
resposta do gestor \cite{brasilMgi2024ppaParticipativo,govuk2021usingData}.

Modelos oficiais e instrumentos consolidados de avaliação de serviços públicos
contribuem para a mensuração da satisfação e da qualidade percebida. Ainda assim,
permanecem desafios relacionados à baixa adesão, ao esforço exigido do usuário e
à transformação das respostas em informação acionável para o gestor. Avaliações
longas podem reduzir a participação; avaliações curtas podem gerar indicadores
pouco explicativos; dados de navegação sem interpretação podem produzir apenas
registros técnicos sem valor gerencial. Esses limites dialogam com discussões
sobre taxa de resposta em pesquisas digitais, uso de dados de desempenho e
sucesso de sistemas de informação
\cite{delightedResponseRate,surveymonkeyResponseRate,govuk2021usingData,deloneMclean2003}.

Este trabalho parte da seguinte questão de pesquisa: como propor e prototipar um
avaliador de serviços públicos digitais que identifique sinais de dificuldade na
navegação, colete feedback contextual de baixo esforço e organize informações
mais úteis para a melhoria dos serviços avaliados do que um formulário estático
exibido apenas ao final da jornada?

\section{Objetivo Geral}

Propor e desenvolver um protótipo inicial de avaliador de serviços públicos
digitais capaz de registrar sinais de dificuldade na navegação, coletar feedback
contextual estruturado, comparar esse fluxo com uma avaliação estática e
apresentar indicadores administrativos que apoiem a melhoria contínua dos
serviços.

\section{Objetivos Específicos}

Para alcançar o objetivo geral, foram definidos os seguintes objetivos
específicos:

\begin{itemize}
    \item levantar fundamentos legais, técnicos e teóricos sobre avaliação de
    serviços públicos, qualidade percebida, usabilidade, governo digital e Lean
    Inception;
    \item definir, por meio de Lean Inception, o fluxo de detecção de
    dificuldade, sugestão de redirecionamento e coleta de feedback;
    \item construir um protótipo funcional com API, páginas simuladas de serviço
    público, feedback contextual, formulário estático comparativo e visualização
    administrativa;
    \item estruturar os dados e indicadores necessários para apoiar a análise de
    cliques, tempo de navegação, taxa de resposta, avaliações, NPS e pontos
    recorrentes de melhoria.
\end{itemize}

A estrutura dos objetivos combina a fundamentação legal e técnica da avaliação de
serviços públicos com práticas de definição de MVP e prototipação de software
\cite{brasil2017lei13460,brasil2021lei14129,caroli2018leanInception,Sommerville2007}.

\section{Justificativa}

Mecanismos digitais de avaliação podem fortalecer a participação do usuário na
melhoria dos serviços públicos. Os materiais da Enap destacam que a avaliação
centrada no cidadão desloca o foco das ferramentas tecnológicas para as
necessidades, demandas e satisfação dos usuários \cite{limaPereira2020modulo1}.
Também indicam que avaliações devem alimentar ciclos de melhoria, permitindo
identificar problemas recorrentes, atributos críticos e prioridades de
intervenção \cite{limaPereira2020modulo4}.

A proposta é relevante porque muitos usuários já conhecem avaliações simples por
estrelas, motivos e comentários. Uma solução pública pode aproveitar essa
familiaridade para reduzir o esforço de resposta, desde que use os dados para
finalidades compatíveis com o interesse público. Além disso, a combinação entre
eventos de navegação e feedback declarado permite que o gestor observe tanto o
que o usuário informou quanto o que ocorreu durante a jornada digital.

Do ponto de vista técnico-científico, o trabalho aproxima qualidade de serviço,
experiência do usuário, design centrado no usuário, sucesso de sistemas de
informação, prototipação e melhoria pública. Essa aproximação é relevante porque
o problema não é apenas criar mais um formulário, mas projetar um artefato capaz
de transformar sinais de navegação e percepções declaradas em dados úteis para
decisão. Na prática, a contribuição ganha forma em uma API, em telas navegáveis,
em regras de detecção de dificuldade e em indicadores administrativos que podem
ser avaliados, criticados e refinados a partir de uma base técnica explícita
\cite{deloneMclean2003,hevner2004designscience,Sommerville2007}.

\section{Escopo Inicial}

O escopo inicial concentra-se na proposta e prototipação de uma solução. O
trabalho não pretende implantar um sistema em ambiente governamental real, nem
substituir instrumentos oficiais de avaliação já existentes. A solução é
complementar e exploratória: simula páginas de serviço público, registra eventos
de navegação, aplica regras configuráveis de detecção de dificuldade, abre uma
coleta simples de feedback, oferece formulário estático para comparação e
apresenta dados administrativos agregados. Essa delimitação segue a lógica de
prototipação como forma de reduzir incertezas antes de uma solução definitiva e
mantém cuidado com minimização de dados pessoais
\cite{Sommerville2007,brasil2018lgpd}.

A validação prevista para esta etapa será preliminar, por análise do protótipo,
demonstração de uso e revisão da coerência entre o fluxo proposto, os fundamentos
teóricos e os dados coletados. Testes com usuários reais, integração com portais
governamentais, calibração estatística das faixas esperadas e avaliação de
impacto ficam como evolução possível para etapas posteriores.
